% ============================================================================
%  Curatio — Full Presentation Walkthrough & Speaker Script
%  Compile with:  pdflatex "presentation-walkthrough.tex"  (run twice for TOC)
%  Needs: amsmath, amssymb, xcolor, tcolorbox, enumitem, booktabs, hyperref
% ============================================================================
\documentclass[11pt]{article}

\usepackage[utf8]{inputenc}
\usepackage[T1]{fontenc}
\usepackage[a4paper,margin=1in]{geometry}
\usepackage{amsmath,amssymb}
\usepackage{xcolor}
\usepackage[most]{tcolorbox}
\usepackage{enumitem}
\usepackage{booktabs}
\usepackage{array}
\usepackage{fancyhdr}
\usepackage{titlesec}
\usepackage[hidelinks]{hyperref}

% ---- Theme colours (match the Curatio app) --------------------------------
\definecolor{forest}{HTML}{166534}
\definecolor{vibrant}{HTML}{22C55E}
\definecolor{deepforest}{HTML}{064E3B}
\definecolor{slate}{HTML}{1E293B}
\definecolor{mutel}{HTML}{475569}
\definecolor{sand}{HTML}{F4F7F5}
\definecolor{satRed}{HTML}{DC143C}
\definecolor{satOrange}{HTML}{FF8C00}
\definecolor{satYellow}{HTML}{D4A500}
\definecolor{satGreen}{HTML}{228B22}
\definecolor{satBlue}{HTML}{1E90FF}

\hypersetup{linkcolor=forest, urlcolor=forest, citecolor=forest}

% ---- Section styling -------------------------------------------------------
\titleformat{\section}{\Large\bfseries\color{forest}}{\thesection}{0.6em}{}
\titleformat{\subsection}{\large\bfseries\color{deepforest}}{\thesubsection}{0.6em}{}
\setlist{nosep, leftmargin=1.4em, topsep=2pt}

% ---- Coloured callout boxes ------------------------------------------------
\newtcolorbox{saybox}{
  enhanced, breakable, colback=forest!6, colframe=forest,
  boxrule=0.8pt, arc=3pt, left=8pt, right=8pt, top=6pt, bottom=6pt,
  fonttitle=\bfseries\color{white}, coltitle=white,
  title={SAY THIS}, attach title to upper={\par\medskip}}

\newtcolorbox{onscreen}{
  enhanced, breakable, colback=slate!4, colframe=slate!55,
  boxrule=0.6pt, arc=3pt, left=8pt, right=8pt, top=6pt, bottom=6pt,
  fonttitle=\bfseries, coltitle=slate,
  title={ON SCREEN}, attach title to upper={\par\smallskip}}

\newtcolorbox{idea}{
  enhanced, breakable, colback=vibrant!7, colframe=vibrant!70!black,
  boxrule=0.6pt, arc=3pt, left=8pt, right=8pt, top=6pt, bottom=6pt,
  fonttitle=\bfseries, coltitle=forest,
  title={THE IDEA (so you understand it)}, attach title to upper={\par\smallskip}}

\newtcolorbox{demo}{
  enhanced, breakable, colback=satOrange!8, colframe=satOrange!85!black,
  boxrule=0.6pt, arc=3pt, left=8pt, right=8pt, top=6pt, bottom=6pt,
  fonttitle=\bfseries, coltitle=satOrange!60!black,
  title={LIVE DEMO --- flip the card}, attach title to upper={\par\smallskip}}

\newtcolorbox{qabox}{
  enhanced, breakable, colback=slate!5, colframe=deepforest,
  boxrule=0.6pt, arc=3pt, left=8pt, right=8pt, top=6pt, bottom=6pt,
  fonttitle=\bfseries, coltitle=deepforest,
  title={LIKELY QUESTIONS}, attach title to upper={\par\smallskip}}

% Inline coloured SATS chips
\newcommand{\chip}[2]{\textcolor{#1}{\textbf{#2}}}
\newcommand{\transition}[1]{\medskip\noindent\textit{\textcolor{mutel}{Transition: #1}}\medskip}
\newcommand{\route}[1]{{\small\texttt{\textcolor{mutel}{#1}}}}

% Slide header helper:  \slide{n}{title}{route}
\newcommand{\slide}[3]{%
  \subsection{Slide #1 --- #2}%
  \noindent\route{Route: #3}\par\smallskip}

\pagestyle{fancy}
\fancyhf{}
\fancyhead[L]{\small\textcolor{forest}{Curatio}}
\fancyhead[R]{\small\textcolor{mutel}{Presentation Walkthrough \& Speaker Script}}
\fancyfoot[C]{\small\thepage}
\renewcommand{\headrulewidth}{0.4pt}

% ============================================================================
\begin{document}

\begin{titlepage}
\centering
\vspace*{3cm}
{\Huge\bfseries\color{forest} Curatio}\\[0.4cm]
{\LARGE\bfseries A Hospital Triage Chatbot for Colour-Coded Clinical Pathways}\\[1cm]
{\Large\itshape Full Presentation Walkthrough \&\ Speaker Script}\\[0.6cm]
{\large What every slide shows, what it means, and \textbf{exactly what to say}.}\\[2cm]
{\large Final Year Project}\\[0.3cm]
{\large \today}\\[2cm]
\begin{tcolorbox}[width=0.85\textwidth, colback=sand, colframe=forest, arc=4pt]
\small\centering
This document mirrors the live web presentation at \texttt{/dashboard}. Each section matches one
slide \emph{in click order}. Read the \textbf{SAY THIS} boxes aloud (or paraphrase) and follow the
\textbf{LIVE DEMO} cues to flip the interactive formula cards on screen.
\end{tcolorbox}
\end{titlepage}

\tableofcontents
\newpage

% ===========================================================================
\section{How to use this script}

\begin{itemize}
  \item \textbf{Structure.} There are \textbf{17 slides} in the live app, in a fixed click order.
        This document has one subsection per slide, in that same order.
  \item \textbf{The coloured boxes.}
    \begin{itemize}
      \item \textbf{ON SCREEN} --- what the audience is looking at, so you can point at it.
      \item \textbf{THE IDEA} --- the concept in plain terms, so \emph{you} fully understand it.
      \item \textbf{SAY THIS} --- a natural spoken script. Read it or make it your own.
      \item \textbf{LIVE DEMO} --- when to click a ``flip / See worked example'' card on the slide.
      \item \textbf{LIKELY QUESTIONS} --- what an examiner may ask, with a ready answer.
    \end{itemize}
  \item \textbf{Golden rule.} You do not need to recite formulas. Your job is to explain
        \emph{what each formula does} and \emph{why it makes the system safer}. The maths is on the
        screen as evidence of rigour.
\end{itemize}

\subsection*{Pre-flight checklist (5 minutes before you start)}
\begin{enumerate}
  \item Start the app and open the home page; click \textbf{Presentation} to land on \route{/dashboard}.
  \item Confirm the formula \emph{flip cards} animate when clicked (test one).
  \item Have the \textbf{Test Model} page (\route{/test}) open in a second tab in case you want a live prediction.
  \item Full-screen the browser; hide the bookmarks bar.
\end{enumerate}

\subsection*{Suggested timing (about 12--15 minutes)}
\begin{itemize}
  \item Foundations (slides 1--2): 2 min.
  \item The 5-step pipeline (slides 3--10): 6--7 min.
  \item Parallel safety branches \&\ deep maths (slides 11--15): 3--4 min.
  \item Worked example + Red protocol + close (slides 16--17): 2 min.
\end{itemize}

% ===========================================================================
\section{The 60-second big picture (memorise this)}

\begin{idea}
Curatio is a \textbf{chatbot that triages patients}. A patient describes their symptoms (by typing or
by voice). The system reads that, optionally reads their vital signs, and outputs a \textbf{single SATS
colour} --- Red, Orange, Yellow, Green or Blue --- that tells the hospital \emph{how fast} the patient
must be seen and \emph{which department} they go to. It is built and validated around
\textbf{Komfo Anokye Teaching Hospital (KATH)} in Kumasi, Ghana.
\end{idea}

\noindent The decision is produced by \textbf{four cooperating layers}, which the formal model writes as one
``session state'':
\[
\mathbf{S}_s = (X,\; \mathbf{D},\; T,\; \mathbf{P},\; C,\; \text{flags})
\]
\[
\mathbf{D} = f_{\mathrm{NLP}}(X),\quad
T = f_{\mathrm{TEWS}}(\mathbf{v}),\quad
\mathbf{P} = f_{\mathrm{Bayes}}(E),\quad
C = f_{\mathrm{fusion}}(\mathbf{D}, T, \mathbf{P})
\]

\begin{itemize}
  \item $X$ --- the patient's complaint as English text.
  \item $\mathbf{D} = f_{\mathrm{NLP}}(X)$ --- \textbf{BioBERT} reads the language and estimates urgency.
  \item $T = f_{\mathrm{TEWS}}(\mathbf{v})$ --- a deterministic \textbf{vital-signs} score (no AI).
  \item $\mathbf{P} = f_{\mathrm{Bayes}}(E)$ --- a \textbf{Bayesian} fallback for when data is missing or uncertain.
  \item $C = f_{\mathrm{fusion}}(\dots)$ --- a \textbf{fusion} step that combines everything into the final colour, always erring on the side of caution.
\end{itemize}

\begin{saybox}
``Curatio is an AI triage assistant. A patient tells our chatbot what's wrong --- in text or in their
own language by voice --- and Curatio turns that into a colour-coded urgency level that routes them to
the right part of the hospital. Under the hood it combines four layers: a fine-tuned medical language
model called BioBERT, a vital-signs score called TEWS, a Bayesian fallback for uncertainty, and a
fusion rule that always picks the \emph{safer} answer when those layers disagree. We built it for the
South African Triage Scale as used at Komfo Anokye Teaching Hospital in Kumasi.''
\end{saybox}

\subsection*{Opening lines (use these to start)}
\begin{saybox}
``Good [morning/afternoon] everyone. Emergency departments are overwhelmed, and the first, most critical
decision is \emph{triage} --- deciding who is seen first. My project, \textbf{Curatio}, is a chatbot that
performs that first triage step automatically and explainably. Over the next few minutes I'll walk you
through exactly how a patient's words become a safe, colour-coded clinical decision. Let's start with
what triage actually is.''
\end{saybox}

% ===========================================================================
\section{Part I --- Foundations}

% ---------------------------------------------------------------------------
\slide{1}{What is Triage?}{/dashboard}

\begin{onscreen}
A lead paragraph defining triage; an \emph{urgency bar}; and a table of the five SATS colours with their
target times and clinical meaning; plus a ``plain English'' traffic-light analogy.
\end{onscreen}

\begin{center}\small
\begin{tabular}{@{}lll@{}}
\toprule
\textbf{Colour} & \textbf{Target time} & \textbf{Clinical meaning}\\
\midrule
\chip{satRed}{Red}       & 0 min      & Immediate resuscitation --- life-threatening\\
\chip{satOrange}{Orange} & $<$ 10 min & Very urgent --- high dependency\\
\chip{satYellow}{Yellow} & $<$ 60 min & Urgent --- needs physician review\\
\chip{satGreen}{Green}   & $<$ 4 hours& Non-urgent --- routine care\\
\chip{satBlue}{Blue}     & $<$ 2 hours& Deceased on arrival --- dignity protocol\\
\bottomrule
\end{tabular}
\end{center}

\begin{idea}
Triage = sorting patients by \emph{urgency}, not by arrival order, so the sickest are seen first. SATS
(South African Triage Scale) is a validated, widely-used colour system. It is the clinical standard our
whole system speaks in.
\end{idea}

\begin{saybox}
``Triage is the process of sorting patients by clinical urgency, so the sickest are seen first --- not
just whoever arrived first. We use the \textbf{South African Triage Scale}, which has been validated at
Komfo Anokye Teaching Hospital. Think of it as a traffic light for the emergency department: each colour
tells staff how quickly a patient must be seen and which area they go to. \chip{satRed}{Red} means
resuscitate \emph{now}; \chip{satOrange}{Orange}, within ten minutes; \chip{satYellow}{Yellow}, within an
hour; \chip{satGreen}{Green}, routine care. Importantly, this is a \emph{decision-support} tool --- it
supports clinical staff, it does not replace their judgment.''
\end{saybox}

\transition{``So those are the colours. The next question is: what does each colour actually \emph{do} inside a Ghanaian hospital?''}

% ---------------------------------------------------------------------------
\slide{2}{Colour $\rightarrow$ Pathways $\rightarrow$ Departments}{/dashboard/pathways}

\begin{onscreen}
A flow diagram from colour to pathway, and a table mapping each colour to a Ghana destination and the
automated ``system action''.
\end{onscreen}

\begin{center}\small
\begin{tabular}{@{}llll@{}}
\toprule
\textbf{Colour} & \textbf{Pathway} & \textbf{Ghana destination} & \textbf{System action}\\
\midrule
\chip{satRed}{Red}       & RED Protocol    & Resuscitation Room       & Code Red; bypass registration \& payment\\
\chip{satOrange}{Orange} & ORANGE Protocol & Majors (ED)              & 10-min countdown; cardiology alert\\
\chip{satYellow}{Yellow} & YELLOW Protocol & Majors (ED)              & Standing orders (e.g.\ malaria test)\\
\chip{satGreen}{Green}   & GREEN Protocol  & Minors / OPD / Polyclinic& Divert non-urgent cases from the ER\\
\chip{satBlue}{Blue}     & BLUE Protocol   & Mortuary                 & Silent dignity protocol; notify social workers\\
\bottomrule
\end{tabular}
\end{center}

\begin{idea}
The colour is not just a label --- it triggers a concrete \emph{workflow}: where the patient physically
goes, and what the system does automatically (alarms, countdowns, diverting non-urgent patients away
from the ER to reduce congestion).
\end{idea}

\begin{saybox}
``Each colour maps to a real clinical pathway and a specific department. A \chip{satRed}{Red} patient
triggers a Code Red and is rushed to the Resuscitation Room --- they even bypass registration and
payment, because seconds matter. \chip{satOrange}{Orange} starts a ten-minute countdown and alerts
cardiology. \chip{satYellow}{Yellow} patients wait in Majors but can have standing orders started, like a
malaria test, so no time is wasted. And critically, \chip{satGreen}{Green} patients are diverted to the
outpatient or polyclinic queue --- this is how Curatio reduces overcrowding in the emergency room.
Curatio performs this routing automatically based on the colour it assigns.''
\end{saybox}

\transition{``Now that you know the outputs, let me show you how a patient actually enters the system.''}

% ===========================================================================
\section{Part II --- The five-step pipeline}

\begin{idea}
Slides 3--10 follow one patient through the pipeline: \textbf{(1) input $\rightarrow$ (2) text example
$\rightarrow$ (3) voice $\rightarrow$ (4) BioBERT $\rightarrow$ (5) fusion}. A progress stepper at the top
of each slide shows where you are. Keep saying ``Step $n$ of 5'' so the audience never gets lost.
\end{idea}

% ---------------------------------------------------------------------------
\slide{3}{How Data Enters the System (Step 1 of 5)}{/dashboard/pipeline/input}

\begin{onscreen}
A ``Step 1 of 5'' badge, the pipeline stepper, a fork diagram (\emph{typed text} vs \emph{voice note}),
and a roadmap of the whole journey.
\end{onscreen}

\begin{idea}
There are two ways a patient can give input --- \textbf{typing} or \textbf{voice} --- but both converge
into the same downstream pipeline. Voice matters because patients in Ghana may speak Twi or another local
language, not English.
\end{idea}

\begin{saybox}
``Every triage session begins with the patient describing their symptoms to our chatbot. They can do this
in two ways: they can \textbf{type} a message --- for example, `I have chest pain' --- or they can send a
\textbf{voice note}, often in Twi or another local Ghanaian language. Two ways in, but one shared pipeline
afterwards. Over the next slides we'll follow that path: a text example, then the speech-and-translation
step, then our BioBERT model, then fusion into a final colour. Vital signs and the Bayesian fallback join
later, at the fusion stage.''
\end{saybox}

\transition{``Let's look at what the simplest case --- typed text --- actually looks like to the system.''}

% ---------------------------------------------------------------------------
\slide{4}{Example Text Input (Step 2 of 5)}{/dashboard/pipeline/text-example}

\begin{onscreen}
A chat mock-up showing the message \texttt{``thunderclap headache, worsening with movement''}; the raw
JSON the system stores; and the nurse's ground-truth label: \chip{satOrange}{Acuity Level 2 --- Highly urgent}.
\end{onscreen}

\begin{idea}
When a patient types, the message is stored verbatim as raw text and becomes the model input we call
\textbf{Text $X$} (capped at 128 tokens). This example is a \emph{real row} from our training data, which
makes the demo credible.
\end{idea}

\begin{saybox}
``If the patient types, we store their message directly as raw text. Here's a real example from our
dataset: `thunderclap headache, worsening with movement.' A `thunderclap' headache is a classic red-flag
phrase for a possible brain bleed, so the nurse ground-truth label here is \chip{satOrange}{Level 2,
highly urgent}. This text becomes what we call \textbf{Text X} --- the input to our model, limited to 128
tokens. If the patient had used voice instead, we'd first need to transcribe and translate it, which is
the very next step.''
\end{saybox}

\transition{``So what happens when the input isn't clean English text, but a voice note in Twi?''}

% ---------------------------------------------------------------------------
\slide{5}{Speech \& Translation (Step 3 of 5)}{/dashboard/pipeline/voice}

\begin{onscreen}
The audio$\rightarrow$translate diagram and a math card for the voice pipeline, with a
\textbf{``See worked example''} flip button. Two ``walkthrough'' chips contrast the typed path vs the voice path.
\end{onscreen}

\[
X = g_{\mathrm{trans}}\!\bigl(g_{\mathrm{ASR}}(\text{audio})\bigr)
\]

\begin{idea}
Read the formula \emph{inside-out}: take the audio, run \textbf{ASR} (automatic speech recognition,
$g_{\mathrm{ASR}}$) to get words, then run \textbf{translation} ($g_{\mathrm{trans}}$) into medical
English. The result is the same $X$ that typed text produces. BioBERT only understands English, so this
step is what makes the system usable for non-English speakers.
\end{idea}

\begin{demo}
Click \textbf{``See worked example''} on the voice formula. It shows patient TG-001 saying a Twi phrase,
$g_{\mathrm{ASR}}$ transcribing it to Twi text, $g_{\mathrm{trans}}$ translating it to ``I have chest pain,
I feel unwell'', and that English string becoming $X$.
\end{demo}

\begin{saybox}
``If the patient sends audio, we transcribe and translate it. Read the formula from the inside out: the
audio goes through $g_{\mathrm{ASR}}$ --- speech-to-text, basically ears writing the words down --- and
then through $g_{\mathrm{trans}}$, a medical interpreter that converts the local language into English.
The output is exactly the same Text $X$ that a typed message produces. [\emph{Flip the card}.] Here a
patient speaks in Twi, we transcribe it, translate it to `I have chest pain, I feel unwell', and that
becomes $X$. The key point: BioBERT only reads English, so everyone's input is normalised to English
before it reaches the model.''
\end{saybox}

\transition{``Now both paths have produced Text X. This is where the intelligence happens --- our fine-tuned BioBERT.''}

% ---------------------------------------------------------------------------
\slide{6}{Fine-Tuned BioBERT Layer (Step 4 of 5)}{/dashboard/pipeline/biobert}

\begin{onscreen}
A BioBERT flow diagram and three flip-cards (the NLP chain, the softmax, and the confidence gate), plus a
table mapping predicted acuity to SATS colour.
\end{onscreen}

\begin{align*}
&X \xrightarrow{\tau} (t_i) \xrightarrow{E} H^{(0)} \xrightarrow{L\times} H^{(L)} \xrightarrow{\text{head}} \hat{\mathbf{y}} \\[4pt]
&\hat{\mathbf{y}} = \mathrm{softmax}(W h_{\mathrm{[CLS]}} + \mathbf{b}),\qquad \hat{y}_c = P(\text{acuity}=c\mid X)\\[4pt]
&\text{confidence} = \max_c \hat{y}_c,\quad \text{if confidence} < \tau \Rightarrow \text{Bayesian fallback}
\end{align*}

\begin{center}\small
\begin{tabular}{@{}lll@{}}
\toprule
\textbf{Predicted acuity} & \textbf{SATS colour} & \textbf{Meaning}\\
\midrule
Level 1--2 & \chip{satRed}{Red} / \chip{satOrange}{Orange} & Highly urgent symptoms\\
Level 3    & \chip{satYellow}{Yellow}                     & Moderate urgency\\
Level 4--5 & \chip{satGreen}{Green}                       & Non-urgent / routine\\
\bottomrule
\end{tabular}
\end{center}

\begin{idea}
BioBERT is a language model \emph{pre-trained on biomedical text} (PubMed), which we then \emph{fine-tuned}
on our triage data. It reads Text $X$ and outputs five probabilities $\hat{\mathbf{y}}$ --- one per acuity
level. The biggest one wins, and its value is the \textbf{confidence}. If confidence is below a threshold
$\tau \approx 0.85$, we don't blindly trust it --- we flag it for the Bayesian fallback.
\end{idea}

\begin{demo}
Flip the three cards in turn: (1) the NLP chain shows ``thunderclap headache'' tokenised and classified to
Level 2 with $\hat{\mathbf{y}} = [0.01, 0.94, 0.03, 0.01, 0.01]$; (2) the softmax card shows those numbers
as percentages (94\% on Level 2); (3) the confidence card shows a high-confidence pass (0.94) vs a
borderline glaucoma case at 0.854 that gets flagged.
\end{demo}

\begin{saybox}
``Now the intelligence. Text $X$ passes through our fine-tuned BioBERT. BioBERT is a language model that
was pre-trained on millions of medical abstracts, so it already understands clinical English. We then
fine-tuned it on our own triage data so it learned \emph{our hospital's phrasing}. It answers one
question: `how urgent does this sound?' [\emph{Flip card 1}] The text is broken into tokens, passed
through twelve layers, and classified --- here `thunderclap headache' comes out as Level 2 with 94\%
probability. [\emph{Flip card 2}] The softmax just turns the raw scores into percentages that add up to
100\%; the highest wins. [\emph{Flip card 3}] And this is the safety valve: we take the confidence ---
the highest probability --- and if it's below about 0.85, we don't trust BioBERT alone; we flag the case
for our Bayesian fallback. That predicted level then maps to a SATS colour, and goes to fusion in Step 5.''
\end{saybox}

\begin{qabox}
\textit{``Why BioBERT and not regular BERT or ChatGPT?''} --- BioBERT is pre-trained on biomedical
literature, so it understands clinical terminology far better than general BERT, and unlike a large
hosted LLM it is small, runs locally/offline, is deterministic, and is auditable --- all important for a
hospital in a low-connectivity setting.
\end{qabox}

\transition{``Let me open the hood for one slide and show what `passes through BioBERT' actually means mathematically.''}

% ---------------------------------------------------------------------------
\slide{7}{BioBERT: Inside the Model}{/dashboard/math/biobert-inference}

\begin{onscreen}
Four math cards for one forward pass: input embedding, self-attention, the transformer layer ($\times 12$),
and the [CLS] summary + classify. Each flips to a worked example on ``crushing chest pain''.
\end{onscreen}

\begin{align*}
&\textbf{Embed: } && E(t_i) = E_{\text{word}}(t_i) + E_{\text{pos}}(i) + E_{\text{seg}}(t_i)\\
&\textbf{Attention: } && \mathrm{Attention}(Q,K,V) = \mathrm{softmax}\!\Bigl(\tfrac{QK^\top}{\sqrt{d_k}}\Bigr)V\\
&\textbf{Summarise: } && h_{\mathrm{[CLS]}} = H^{(L)}_{1,:}\in\mathbb{R}^{768}
\end{align*}

\begin{idea}
This slide is for the technically curious examiner. The story in four beats: \textbf{(1) Embedding} turns
each word into a vector that also encodes its position; \textbf{(2) Self-attention} lets the model weigh
which words matter --- clinical words like ``crushing'' and ``chest'' get high weight, filler words like
``the'' get almost none; \textbf{(3)} this repeats through \textbf{12 transformer layers}; \textbf{(4)} a
special \texttt{[CLS]} token collects a 768-number summary of the whole sentence, which the classifier
reads to produce the urgency.
\end{idea}

\begin{demo}
Flip the \textbf{self-attention} card --- it shows that for ``crushing chest pain'', attention weights are
\texttt{crushing}=0.38, \texttt{chest}=0.35, \texttt{pain}=0.22, and filler words $<0.05$. This is the
most intuitive one to show.
\end{demo}

\begin{saybox}
``For those who want the detail, here's what actually happens inside. First, each word becomes a vector
that also remembers its position in the sentence. Then comes the key mechanism --- \textbf{self-attention}
--- where the model decides which words matter most. [\emph{Flip the attention card}.] Look: for `crushing
chest pain', the model puts almost all its weight on the clinical words --- crushing, chest, pain --- and
basically ignores filler words. That happens twelve times, in twelve layers, each one building a deeper
understanding. Finally, a special summary token called \texttt{[CLS]} compresses the whole sentence into
768 numbers, and that summary --- not the individual words --- is what the classifier reads to decide
urgency. So it's really one forward pass: read the sentence, summarise it, output an urgency percentage.''
\end{saybox}

\transition{``That's how it thinks. The natural next question is: how did it \emph{learn} to think that way?''}

% ---------------------------------------------------------------------------
\slide{8}{BioBERT: How It Was Trained}{/dashboard/math/biobert-training}

\begin{onscreen}
Two loss cards (pre-training MLM loss, fine-tuning cross-entropy loss), a hyper-parameter table, a 4-box
training-loop diagram, and a note that vitals were excluded.
\end{onscreen}

\[
\mathcal{L}_{\mathrm{MLM}} = -\!\!\sum_{i\in\mathcal{M}}\!\log P(t_i\mid H^{(L)};\theta)
\qquad\qquad
\mathcal{L} = -\sum_{i=1}^{N}\log P(y_i\mid X_i;\theta)
\]

\begin{center}\small
\begin{tabular}{@{}ll@{}}
\toprule
\textbf{Hyper-parameter} & \textbf{Value}\\
\midrule
Base model    & \texttt{Yuvrajxms09/biobert-triage-classifier}\\
Merge         & \texttt{train.csv} $\bowtie$ \texttt{chief\_complaints.csv} ($\sim$80k rows)\\
Epochs        & 3 (13{,}500 steps)\\
Best \texttt{eval\_loss} & 0.001812\\
\bottomrule
\end{tabular}
\end{center}

\begin{idea}
Two stages. \textbf{Pre-training} is ``fill-in-the-blank'' on PubMed --- the MLM loss rewards the model
for guessing a masked word like ``chest'' from context; this is how it learned medical English in general.
\textbf{Fine-tuning} is our part: we minimise the cross-entropy loss $\mathcal{L}$ between the model's
predicted acuity and the \emph{nurse's} label across 80k complaints. The training loop is just:
\emph{forward $\rightarrow$ measure loss $\rightarrow$ back-propagate $\rightarrow$ repeat} for 3 epochs.
We deliberately trained on \emph{language only}, no vitals.
\end{idea}

\begin{demo}
Flip the fine-tuning card: one row --- ``contraception advice, intermittent'', nurse label Level 5 ---
the model predicts Level 5 at 97\%, so the loss it adds is $-\log(0.97)\approx0.03$, a tiny error.
\end{demo}

\begin{saybox}
``How did BioBERT learn? In two stages. First, \textbf{pre-training}: before we ever touched it, it was
trained to fill in blanks in medical sentences --- given `the patient presented with severe blank pain',
predict `chest'. After millions of these, it understands medical language. Then \textbf{we fine-tuned} it
on our task: for 80,000 real chief complaints, it guessed an acuity level, we compared it to the nurse's
label, and we nudged the weights to reduce the error. [\emph{Flip the card}.] For instance, for
`contraception advice' the nurse said Level 5, the model said Level 5 at 97\% confidence, so the error is
tiny. We ran this for three epochs and the validation loss dropped to about 0.0018. One deliberate choice:
we trained on \emph{language alone} --- no vital signs --- so the model learns urgency purely from how the
complaint is phrased. The vitals come in separately, which keeps the two signals independent.''
\end{saybox}

\transition{``So how well does it actually do at scale? Let me show you the data and the findings.''}

% ---------------------------------------------------------------------------
\slide{9}{Data Exploration \& Findings}{/dashboard/math/data-exploration}

\begin{onscreen}
Three tables --- dataset sizes, the acuity distribution of the training set, and example predictions ---
plus a ``key findings'' list.
\end{onscreen}

\begin{center}\small
\begin{tabular}{@{}lll@{}}
\toprule
\textbf{Dataset} & \textbf{Rows} & \textbf{Role}\\
\midrule
\texttt{train.csv} & 80{,}000 & Labels + vitals (80k acuity labels)\\
\texttt{test.csv} & 20{,}000 & Held-out evaluation (no labels)\\
\texttt{chief\_complaints.csv} & 100{,}000 & Raw symptom text input\\
\texttt{triage\_predictions\_results.csv} & 100{,}000 & Model inference output\\
\bottomrule
\end{tabular}
\quad
\begin{tabular}{@{}lrr@{}}
\toprule
\textbf{Acuity} & \textbf{Count} & \textbf{\%}\\
\midrule
Level 1 & 3{,}222  & 4.0\%\\
Level 2 & 13{,}439 & 16.8\%\\
Level 3 & 28{,}921 & 36.2\%\\
Level 4 & 23{,}020 & 28.8\%\\
Level 5 & 11{,}398 & 14.2\%\\
\bottomrule
\end{tabular}
\end{center}

\begin{idea}
We worked with about 100k records. The acuity distribution is realistic --- most patients are
mid-urgency (Level 3), few are critical (Level 1). Two headline findings: \textbf{99.6\%} of predictions
come out at full confidence (1.0), and \textbf{361} rows sit in the uncertain 0.6--0.99 band --- those are
exactly the cases the Bayesian fallback exists for. Be honest about the limitation: a model this confident
can occasionally be \emph{confidently wrong}, which is precisely why fusion with TEWS and Bayesian exists.
\end{idea}

\begin{saybox}
``We trained and evaluated on roughly 100,000 patient records. The distribution is clinically realistic
--- about a third are Level 3, moderate urgency, and only 4\% are the most critical Level 1. Two findings
stand out. First, 99.6\% of our predictions come out at full confidence. Second --- and more interesting
--- 361 rows land in a borderline confidence band; for example, two glaucoma complaints at around 0.85.
Those are the cases we hand to the Bayesian fallback rather than trusting blindly. And I want to be honest
about the limitation: a very confident model can sometimes be confidently \emph{wrong}, especially on
rare phrasings. That's the whole reason the next layers --- vitals and Bayesian --- exist: to catch what
the language model alone might miss.''
\end{saybox}

\transition{``Which brings us to the heart of the system: fusion --- how we combine the language model with everything else.''}

% ---------------------------------------------------------------------------
\slide{10}{Fusion $\rightarrow$ Triage Colour (Step 5 of 5)}{/dashboard/pipeline/fusion}

\begin{onscreen}
The fusion pipeline diagram, the master fusion-function card (flippable), and an example output chip
\chip{satOrange}{C = Orange $\rightarrow$ Majors (ED)}.
\end{onscreen}

\[
C = f_{\mathrm{fusion}}(\mathbf{D},\, T,\, \mathbf{P},\; \text{flags})
\]

\begin{idea}
Fusion is the master controller. The single most important sentence on this slide: \textbf{fusion is a
safety checklist, not an average}. When the layers disagree, it always takes the \emph{more urgent}
signal. So a strong language signal can override mild vitals --- the system fails safe.
\end{idea}

\begin{demo}
Flip the fusion card: a chest-pain patient where BioBERT says Orange (94\%), TEWS says only Yellow ($T=4$),
and Bayesian says Orange ($P=0.89$) --- fusion outputs \chip{satOrange}{Orange}, because symptoms beat
mild vitals.
\end{demo}

\begin{saybox}
``This is the heart of the system. Fusion takes the three signals --- BioBERT's language reading, the
TEWS vitals score, and the Bayesian estimate --- and combines them into one final colour $C$. The crucial
thing to understand is that fusion is \textbf{not an average} --- it's a \textbf{safety checklist}. When
the layers disagree, it always picks the \emph{more urgent} answer. [\emph{Flip the card}.] Here's a
chest-pain patient: the language model says Orange, but the vitals alone only say Yellow. A naive average
would split the difference and under-triage them. Fusion instead enforces Orange --- the safer call ---
and sends them to Majors with a cardiology alert. The next few slides show the two parallel branches that
feed this fusion step: the vitals score, and the Bayesian fallback.''
\end{saybox}

\transition{``Let's look at the first parallel branch --- turning vital signs into a number.''}

% ===========================================================================
\section{Part III --- The parallel safety branches \& the deep maths}

% ---------------------------------------------------------------------------
\slide{11}{Parallel Path: Vitals $\rightarrow$ TEWS}{/dashboard/pipeline/tews}

\begin{onscreen}
Three flip-cards (TEWS summation, colour-from-vitals lookup, a numeric example) and a band table mapping
$T$ to a colour.
\end{onscreen}

\[
T = f_{\mathrm{TEWS}}(\mathbf{v}) = \sum_{k=1}^{6} w_k\, f_k(v_k),\qquad w_k = 1
\qquad\qquad
C_{\mathrm{TEWS}}(T)=\begin{cases}
\text{Red} & T>7\\ \text{Orange} & 5\le T\le 6\\ \text{Yellow} & 3\le T\le 4\\ \text{Green} & T\le 2
\end{cases}
\]

\begin{idea}
TEWS = Triage Early Warning Score. It's \textbf{pure arithmetic, no AI}: six vital signs (heart rate,
respiratory rate, mobility, temperature, consciousness/AVPU, trauma) each score points; you add them up to
get $T$; $T$ maps to a colour band. Because it's deterministic, it's fully auditable --- every point
traces to a clinical table.
\end{idea}

\begin{demo}
Flip the cards for our running patient: HR\,=\,125 scores 2, RR\,=\,26 scores 2, everything else 0, so
$T=4$ $\rightarrow$ \chip{satYellow}{Yellow} from vitals alone.
\end{demo}

\begin{saybox}
``The first parallel branch is vital signs. We use TEWS --- the Triage Early Warning Score. Unlike
BioBERT, there is \textbf{no AI here at all}; it's deliberate, transparent arithmetic. Six vitals --- heart
rate, breathing rate, mobility, temperature, consciousness, and trauma --- each contribute points, and we
add them into a total $T$. [\emph{Flip the card}.] For our patient, a heart rate of 125 scores 2 points, a
breathing rate of 26 scores another 2, the rest are normal, so $T=4$. We then look $T$ up in the band
table: 3 to 4 means \chip{satYellow}{Yellow}. The reason we keep this deterministic is \emph{auditability}
--- a clinician can check every single point against a standard table.''
\end{saybox}

\transition{``If an examiner wants the full vitals mathematics, the next slide has every piecewise function.''}

% ---------------------------------------------------------------------------
\slide{12}{TEWS: Full Mathematics}{/dashboard/math/tews}

\begin{onscreen}
Six flip-cards: the total, the heart-rate function $f_1$, the respiratory function $f_2$,
score-to-colour, a worked example, and the ``missing vitals'' partial score. Plus a band-to-Ghana-routing
table.
\end{onscreen}

\[
f_1(\mathrm{HR})=\begin{cases}3 & \mathrm{HR}\ge130\\ 2 & 111\le \mathrm{HR}\le129\\ 0 & 51\le \mathrm{HR}\le100\\ 2 & \mathrm{HR}\le40\\ 1 & \text{otherwise}\end{cases}
\qquad
T_{\mathrm{partial}}=\!\!\sum_{k\in\mathcal{K}_{\mathrm{obs}}}\!\! w_k f_k(v_k)
\]

\begin{idea}
This is the ``show your rigour'' slide. The piecewise functions are just clinical lookup tables written
formally. The genuinely important one is \textbf{$T_{\mathrm{partial}}$}: if some vitals are missing
(common in a busy ER), we only sum the \emph{observed} ones and set a \texttt{tews\_incomplete} flag ---
and on incomplete data the system refuses to assign a reassuring Green, deferring instead to the Bayesian
estimate.
\end{idea}

\begin{demo}
Flip $f_1$: for HR\,=\,125, check top-down --- $\ge130$? no; $111\!-\!129$? yes $\rightarrow$ 2 points.
Then flip ``Missing vitals'': with only HR and RR measured, $T_{\mathrm{partial}}=4$ and the system will
not assign Green on partial data.
\end{demo}

\begin{saybox}
``For completeness, here's the full TEWS mathematics. Each vital has a piecewise function --- which is
really just the official clinical lookup table written as maths. [\emph{Flip $f_1$}.] For a heart rate of
125, you read top to bottom: is it 130 or more? No. Is it between 111 and 129? Yes --- so 2 points. The one
I want to highlight is \textbf{partial scoring}. In a real, busy emergency department, you often don't have
every vital. So we only add up the vitals we actually measured, and we raise an `incomplete' flag. And our
safety rule says: on incomplete data, we will \emph{not} hand out a reassuring Green --- instead we defer
to the Bayesian estimate. Every score here is traceable to a SATS table row, so it's auditable for clinical
review.''
\end{saybox}

\transition{``That deferral I keep mentioning --- to Bayesian --- is the second parallel branch. Let's see it.''}

% ---------------------------------------------------------------------------
\slide{13}{Parallel Path: Bayesian Fallback}{/dashboard/pipeline/bayesian-fallback}

\begin{onscreen}
Three flip-cards --- Bayes' rule for the colour posterior, the argmax decision, and the disease-level
override --- plus a ``when does Bayesian run?'' list.
\end{onscreen}

\[
P(C_k\mid E)=\frac{P(E\mid C_k)\,P(C_k)}{\sum_j P(E\mid C_j)\,P(C_j)}
\qquad
C_{\mathrm{Bayes}}=\arg\max_k P(C_k\mid E)
\]

\begin{idea}
Bayes' rule updates a \emph{prior} belief using \emph{evidence}. In plain terms: start with how common
each colour is at this hospital (the prior $P(C_k)$), then adjust by how well the patient's evidence $E$
fits each colour (the likelihood $P(E\mid C_k)$), and normalise. The colour with the highest resulting
probability wins. There's also a \textbf{disease override}: if a specific dangerous disease (e.g.\ heart
attack) is strongly suspected, upgrade to Red/Orange even if vitals look calm.
\end{idea}

\begin{demo}
Flip the posterior card: for chest pain with partial vitals, priors of 5/15/30/50\% become posteriors of
8/89/2/1\% --- \chip{satOrange}{Orange} wins at 89\%. Then flip the override card: $P(\text{MI}\mid
\text{crushing chest pain})\approx0.82 > \tau_B = 0.75$ $\rightarrow$ upgrade to Red/Orange.
\end{demo}

\begin{saybox}
``The second branch is our Bayesian fallback --- this is what runs when BioBERT is uncertain or vitals are
missing. Bayes' rule is just a principled way of updating a belief with evidence. We start with a prior ---
how common each colour normally is at this hospital --- and we update it using the patient's evidence: the
symptoms and whatever vitals we have. [\emph{Flip the card}.] For a chest-pain patient with only partial
vitals, Green started as the most likely at 50\%, but once we factor in the evidence, \chip{satOrange}{Orange}
jumps to 89\% and wins. We pick the highest --- that's the argmax. And there's a critical safety feature:
a \textbf{disease-level override}. [\emph{Flip the override card}.] If the symptoms strongly suggest
something dangerous like a heart attack --- here an 82\% probability --- we upgrade to Red or Orange even if
the vital signs look mild. That's the layer that catches the `looks fine on paper but is actually
critical' patient.''
\end{saybox}

\transition{``The deep-maths version of that branch is next, if you want the full evidence model.''}

% ---------------------------------------------------------------------------
\slide{14}{Bayesian: Full Mathematics}{/dashboard/math/bayesian}

\begin{onscreen}
Flip-cards for Bayes' rule, the \emph{evidence vector} $E$, the argmax, and the disease override pair.
\end{onscreen}

\[
E=\bigl[\text{entity}_1,\dots,\text{entity}_r,\; v_{k_1},\dots,v_{k_s}\bigr]
\qquad
P(D\mid S)=\frac{P(S\mid D)\,P(D)}{P(S)}
\]

\begin{idea}
The new piece here is \textbf{what evidence $E$ actually is}: it's the medical entities BioBERT's NLP
extracted from the text (e.g.\ \texttt{chest\_pain}, \texttt{radiating\_pain}) \emph{plus} whatever vitals
were measured (e.g.\ HR\,=\,125). So Bayesian doesn't ignore the language --- it reuses the extracted facts.
$P(D\mid S)$ is the same Bayes' rule applied at the \emph{disease} level rather than the colour level.
\end{idea}

\begin{demo}
Flip the ``evidence vector'' card: $E=[\texttt{chest\_pain}, \texttt{radiating\_pain}, \mathrm{HR}{=}125,
\mathrm{RR}{=}26]$ --- showing language and vitals combined into one evidence list.
\end{demo}

\begin{saybox}
``Here's the full Bayesian model. The one extra concept worth pointing out is the \textbf{evidence vector
$E$}. [\emph{Flip the evidence card}.] It combines two sources: the medical entities our NLP pulled out of
the patient's words --- like `chest pain' and `radiating pain' --- together with the vital signs we
measured. So the Bayesian layer isn't a separate guess in the dark; it reuses everything the language
model already understood, plus the vitals. The prior is our first guess from hospital history; the
likelihood is how well the symptoms fit each colour; Bayes combines them into an updated probability. And
the disease override is simply the same rule applied to specific diseases instead of colours.''
\end{saybox}

\transition{``Now we can show the full fusion algorithm that ties all three layers together.''}

% ---------------------------------------------------------------------------
\slide{15}{Fusion: Full Mathematics}{/dashboard/math/fusion}

\begin{onscreen}
The fusion diagram, three flip-cards (the fusion function, the urgency ordering, the safety rule), the
``Algorithm 1'' checklist, and a conflict-resolution table.
\end{onscreen}

\[
\mathrm{ord}(\text{Red})=4 > \mathrm{ord}(\text{Orange})=3 > \mathrm{ord}(\text{Yellow})=2 > \mathrm{ord}(\text{Green})=1
\]
\[
\mathrm{ord}(C)\;\ge\;\max\bigl\{\mathrm{ord}(C_{\mathrm{disc}}),\,\mathrm{ord}(C_{\mathrm{TEWS}}),\,\mathrm{ord}(C_{\mathrm{Bayes}})\bigr\}
\]

\begin{idea}
This formalises ``fail safe''. We give each colour a rank (Red\,=\,4 down to Green\,=\,1). The
\textbf{safety rule} says the final colour's rank must be \emph{at least} the highest rank any layer
suggested --- so the patient can never be downgraded below the most urgent opinion. The
\textbf{Algorithm 1} checklist is the order of decisions: life-threatening words $\rightarrow$ Red; very
urgent words $\rightarrow$ Orange; then the TEWS bands; then missing-vitals $\rightarrow$ Bayesian argmax;
then disease override; else Green.
\end{idea}

\begin{demo}
Flip ``urgency ordering'': TEWS says Yellow (ord 2), language says Orange (ord 3); $\max(2,3)=3$
$\rightarrow$ final \chip{satOrange}{Orange}. The higher rank always wins.
\end{demo}

\begin{saybox}
``This slide makes `fail safe' precise. We rank the colours --- Red is 4, down to Green is 1. The
\textbf{safety rule} is one line but it's the ethical core of the whole project: the final colour must be
\emph{at least as urgent} as the most urgent thing any layer detected. [\emph{Flip the ordering card}.] So
if vitals say Yellow but the language says Orange, we take the maximum --- Orange. A patient is never
downgraded into a less urgent queue. The checklist on the right is the exact decision order: check for
life-threatening language first, then very-urgent language, then the vitals bands, then --- if vitals are
missing --- the Bayesian estimate, then the disease override, and only if nothing fires do we assign
Green. It's a transparent, ordered set of rules, not a black box.''
\end{saybox}

\transition{``Let's put every layer together on one real patient, start to finish.''}

% ===========================================================================
\section{Part IV --- Bringing it together}

% ---------------------------------------------------------------------------
\slide{16}{End-to-End Worked Example}{/dashboard/math/worked-example}

\begin{onscreen}
Three layer cards (BioBERT / TEWS / Bayesian), two flip math cards, and a highlighted result banner:
\chip{satOrange}{ORANGE --- Majors (ED)} with cardiology alert, 10-minute countdown, bed in Majors.
\end{onscreen}

\[
\underbrace{d_{\text{chest}}>\tau}_{\text{language}} + \underbrace{T=4}_{\text{TEWS alone = Yellow}} + \underbrace{P(\text{Orange}\mid E)\approx0.89}_{\text{Bayesian}} \;\Rightarrow\; C=\text{Orange}
\]

\begin{idea}
This is the payoff slide --- it shows the layers \emph{disagreeing} and fusion resolving them safely.
Language says Orange, vitals alone say only Yellow, Bayesian says Orange at 89\%. Fusion takes the safer
Orange. This single example proves the central thesis: combining layers beats any one layer alone.
\end{idea}

\begin{demo}
Flip the ``Fusion resolution'' card to walk the audience through each term of the equation above, then
land on the Orange result banner.
\end{demo}

\begin{saybox}
``Let me bring it all together with one real Kumasi patient: central crushing chest pain. Watch the three
layers. BioBERT reads the language and says \chip{satOrange}{Orange} --- this sounds very urgent. But the
vital signs alone only score $T=4$, which is \chip{satYellow}{Yellow} --- on paper they look okay-ish. The
Bayesian layer, using both the symptoms and the partial vitals, says Orange at 89\%. [\emph{Flip the
card}.] So two layers say Orange, one says Yellow. A simple average would under-triage this patient ---
and chest pain is exactly where you can't afford that. Fusion enforces the safer answer:
\chip{satOrange}{Orange}. That triggers a cardiology alert, a ten-minute countdown, and a bed in Majors.
\emph{This} is the whole point of the project --- no single layer is trusted alone; together they fail
safe.''
\end{saybox}

\transition{``Finally, let me show what happens at the most extreme end --- a Red, life-threatening case.''}

% ---------------------------------------------------------------------------
\slide{17}{Resuscitation (Red)}{/dashboard/pathways/resus}

\begin{onscreen}
A red alert banner --- ``TARGET RESPONSE TIME: 0 MINUTES (IMMEDIATE)'' --- and a numbered 3-step
emergency protocol.
\end{onscreen}

\begin{enumerate}
  \item \textbf{Secure Airway \& Vitals} --- assess airway; attach continuous cardiac, SpO$_2$, and NIBP monitoring.
  \item \textbf{Establish Access} --- two large-bore IVs or IO access; draw stat labs (VBG, lactate, cardiac enzymes).
  \item \textbf{Targeted Intervention} --- ACLS/ATLS protocols by cause (defibrillation, epinephrine, massive transfusion).
\end{enumerate}

\begin{idea}
This closes the loop: a colour isn't abstract --- \chip{satRed}{Red} launches a concrete, time-zero
clinical protocol. It shows the system is grounded in real emergency-medicine workflow, not just an ML
score.
\end{idea}

\begin{saybox}
``And to close the loop --- this is what a \chip{satRed}{Red} actually launches. Target response time: zero
minutes, immediate. The protocol is concrete: first secure the airway and get the patient on continuous
monitoring; second, establish two large-bore IV lines and draw the critical labs; third, begin the
targeted intervention --- ACLS or ATLS depending on the cause. The point I want to leave you with is that
Curatio's output isn't an abstract number --- every colour maps to a real, standardised clinical action.''
\end{saybox}

\subsection*{Closing statement}
\begin{saybox}
``To summarise: Curatio takes a patient's own words --- typed or spoken in their language --- and turns
them into a safe, explainable, colour-coded triage decision. It does this by combining a fine-tuned medical
language model, a transparent vital-signs score, and a Bayesian safety net, all fused by a rule that always
errs toward the safer outcome. It's designed for a real hospital --- Komfo Anokye in Kumasi --- to help
clinicians see the sickest patients first and divert non-urgent cases away from a crowded emergency room.
Thank you --- I'm happy to take questions.''
\end{saybox}

% ===========================================================================
\section{Appendix A --- Formula cheat-sheet}

\begin{tabular}{@{}p{0.30\textwidth}p{0.64\textwidth}@{}}
\toprule
\textbf{Formula} & \textbf{Say it as\dots}\\
\midrule
$\mathbf{S}_s=(X,\mathbf{D},T,\mathbf{P},C,\text{flags})$ & ``Everything we know about this session in one bundle.''\\
$X=g_{\mathrm{trans}}(g_{\mathrm{ASR}}(\text{audio}))$ & ``Transcribe the audio, then translate it to English.''\\
$\hat{\mathbf{y}}=\mathrm{softmax}(Wh_{\mathrm{[CLS]}}+\mathbf{b})$ & ``Turn the model's scores into five urgency percentages.''\\
$\text{confidence}=\max_c\hat{y}_c$ & ``How sure the model is; if it's low, call the fallback.''\\
$T=\sum_{k=1}^6 w_k f_k(v_k)$ & ``Add up points from six vital signs.''\\
$P(C_k\mid E)=\frac{P(E\mid C_k)P(C_k)}{\sum_j P(E\mid C_j)P(C_j)}$ & ``Update our guess about the colour using the evidence.''\\
$\mathrm{ord}(C)\ge\max\{\dots\}$ & ``The final colour is never less urgent than any layer said.''\\
\bottomrule
\end{tabular}

% ===========================================================================
\section{Appendix B --- Question bank (rehearse these)}

\begin{qabox}
\textbf{Q: Is this replacing nurses?} \\
No --- it's decision \emph{support}. It standardises and speeds up the first triage step; the clinician
always remains in control and can override.

\medskip
\textbf{Q: What if the AI is wrong?} \\
That's exactly why there are three independent layers and a fail-safe fusion rule. The system is built so
that an error in one layer (say an over-confident BioBERT) is caught by the vitals score or the Bayesian
override, and fusion always defaults to the \emph{more} urgent colour.

\medskip
\textbf{Q: Why is the vitals score not AI?} \\
Deliberately. TEWS must be transparent and auditable --- a clinician can verify every point against a
standard table. Mixing a deterministic clinical score with the learned language model gives us the best of
both: rigour and nuance.

\medskip
\textbf{Q: How does it handle local languages?} \\
The voice pipeline transcribes speech (ASR) and translates it (e.g.\ Twi $\rightarrow$ English) before the
model sees it, so non-English speakers are first-class users.

\medskip
\textbf{Q: How big is the model / can it run offline?} \\
It's a fine-tuned BioBERT (\textasciitilde110M parameters, a \textasciitilde414\,MB weights file). It runs
on CPU and can be deployed locally --- important for low-connectivity hospital settings.

\medskip
\textbf{Q: What were the main limitations?} \\
Confident misclassification on rare phrasings, dependence on translation quality for voice, and the need
for real-world clinical validation. Fusion with TEWS/Bayesian mitigates the first; the latter two are
future work.
\end{qabox}

\end{document}
